% ============================================================================ DISCUSSION
\section{Discussion}\label{sec:discussion}

\subsection{Design rules from the coupled solution}
Each rule below follows from a result of Sections~\ref{sec:evaluate}--\ref{sec:explore} and is
invisible to the single-domain analysis that would normally make the decision.

\emph{Rank integrations on the converged electrothermal state.} The loss a regulator dissipates
depends on the temperature its integration imposes on it. At room temperature the
substrate-backside regulator is ahead of the embedded one by \DeltaPrtZtwoSzero~W and the
board-backside regulator behind it by \DeltaPrtAbsZthreeSzero~W; on the converged field the order
of each pair reverses, by \LossGapZtwoSzero~W and \LossGapZthreeSzero~W, because the converter's
thermal path changes its conversion loss by more than the routing changes. A placement study
that holds the converter loss fixed ranks the integrations in the wrong order.

\emph{An embedded regulator needs local input decoupling.} The package 48-V network of an
embedded regulator carries whatever waveform the cells draw. Without local input capacitance
the package loss rises by \RiseSoneZone~W and the critical joint leaves the calibrated range of
the solder law; with the input pulsed at the overall ratio there is no admissible operating
point. Directly fed backside regulators are insensitive to this choice.

\emph{Order the tiers for the interconnect technology.} With build-up microvias, the high-side
device, the flying capacitors, the low-side device and the inductors should be stacked in that
order: the reference order needs a via of aspect ratio 5.3, the repartitioned one 0.83, with the
same components, area and height. % TODO(data): + its electrothermal effect (z1_tier_reference)

\emph{A backside regulator needs a feed specification.} The direct 48-V feed is outside the
package and must be specified as a loop resistance below $R^*$; where the backside regulator is
behind at zero feed resistance, no connector recovers the difference.
% TODO(data): R* values

\emph{Size the package interface for the rail it carries.} Integrating the regulator on the
board backside removes the board opening, but sends the kiloampere 1-V current through the ball
field: at 1.5~A per ball it occupies 1334 of the 1369 sites of a 0.8-mm array, leaving almost no
signal balls, and only \LandBallsInside{} balls per rail lie under the cells. Redistributing
the rest laterally in the package land side costs \LandPmin--\LandPmax~W and
\LandDVmin--\LandDVmax~mV for \LandTmax--\LandTmin~\textmu m of copper. Keeping the 1-V
current under the cells requires a ball pitch of \PitchMaxZthree~mm or less.

\emph{Do not use temperature as a proxy for joint response.} The package-backside regulator
runs \dTmaxAbsZtwoSzero~K hotter at its peak than the embedded one, yet its critical joint is
\dTjointAbsZtwoSzero~K cooler, because the ball ring sits away from the converter; the joint
still moves \DuRatioZtwo{} times as far, because its driver is set by the bending of the plate the
balls support rather than by their temperature. A screen on joint temperature would rank the two
the wrong way round.

\subsection{Scope and limitations}
The framework is verified stage by stage against matched references, not against a measured
system; no fabricated DVPD package of this kind exists to measure. The converter loss model
omits core, alternating-current winding, dielectric and dead-time losses, which bias the
absolute loss downward by the same mechanism in every integration. The local solder response is
a per-cycle inelastic strain increment, not a life: no damage law or allowable is applied, and
the constitutive law is calibrated to 125\,$^\circ$C, so responses at hotter joints are flagged
as extrapolations. The mechanics is one-way coupled, which the small feedback of thermal strain on
resistance justifies in the intact regime but not for a damaged or delaminated interface. The board is a 30-mm coupon with a film beneath;
its lateral spreading is bracketed by the board-boundary sensitivity. The Z3 ledger omits the
board-plane and land-side 1-V redistribution and is therefore a lower bound, which the sheet
estimate above brackets. Every result belongs to one specification; the conclusions that
generalize are the mechanisms and rules, not the numbers.
